% Generated draft: requires mathematical and source review.
\documentclass[11pt]{article}
\usepackage[margin=25mm]{geometry}
\usepackage{fontspec}
\setmainfont{Latin Modern Roman}
\usepackage{amsmath,amssymb,mathtools,unicode-math}
\setmathfont{Latin Modern Math}
\usepackage{enumitem,needspace,xurl,hyperref}
\setlength{\parindent}{0pt}
\setlength{\parskip}{5pt}
\setlength{\emergencystretch}{3em}
\newcommand{\N}{\mathbb{N}}
\newcommand{\Z}{\mathbb{Z}}
\newcommand{\Q}{\mathbb{Q}}
\newcommand{\R}{\mathbb{R}}
\newcommand{\C}{\mathbb{C}}
\newcommand{\F}{\mathbb{F}}
\newcommand{\A}{\mathbb{A}}
\newcommand{\PP}{\mathbb{P}}
\newcommand{\E}{\mathbb{E}}
\newcommand{\eps}{\varepsilon}
\newcommand{\dd}{\,\mathrm{d}}
\begin{document}
\section*{UnsolvedMath ID: 10000021}
\textbf{Source problem code:} catalogue entry
\subsection{Definitions and mathematical statement}
Give \(\Z^k\) its nearest-neighbor graph structure. Benjamini asks whether, for \emph{some} finite integer \(d\geq3\), there is a probability measure \(\mu\) on embeddings of \(\Z^2\) into \(\Z^d\) with an analogue of exponential intersection tails.

For two independent samples \(\phi,\psi\sim\mu\), the ordinary counting form of an exponential intersection-tail condition would be
\[
 \Pr(I(\phi,\psi)\geq n)\leq C e^{-cn}\qquad(n\geq1)
\]
for constants \(C<\infty\), \(c>0\), and a specified intersection statistic \(I\). For example, one could count common image vertices or common image edges, but these are different conditions and neither is identified by the excerpt as the intended two-dimensional analogue. The display explains the usual tail terminology rather than selecting an unsupplied definition.

The embedding notion is likewise unspecified: an injective nearest-neighbor graph map, a topological embedding using edge paths, and a metric embedding are different possibilities. Root normalization and how the two embeddings are placed relative to each other also matter. The source does not give these conventions or any invariance requirement on \(\mu\). It therefore asks for an appropriate analogue and a measure realizing it, not an already defined claim for every dimension or an ordinary one-dimensional path measure.
\subsection{Short English statement}
In at least one finite dimension \(d\geq3\), can random embeddings of the whole square lattice have exponentially decaying intersection tails? The embedding convention and the intersection statistic for the requested analogue must be specified.
\subsection{Sources}
\begin{itemize}
\item Local UnsolvedMath registry, permanent ID 10000021, source code catalogue entry; Itai Benjamini; no proposal year supplied; updated 2026-08-06.
\item Benjamini, \emph{Coarse Geometry and Randomness}, Saint-Flour notes, Open problem 4.27, PDF page 37: \url{https://www.wisdom.weizmann.ac.il/~itai/stflouraug24.pdf#page=37}. Supplied archive: \url{https://arquivo.pt/noFrame/replay/20201231041548id_/http://www.wisdom.weizmann.ac.il/~itai/stflouraug24.pdf}. Registry access date 2026-07-29; text extraction checked against rendered PDF, as reported by provenance. Item retained after cross-document deduplication and source-side review; no complete resolution found in the bounded audit. Current status and release reuse terms need review.
\item The summary associates the question with Pemantle--Peres theory but supplies no exact reference or definition of the higher-dimensional intersection-tail property.
\end{itemize}
\end{document}
